\documentclass[11pt,letterpaper]{article}
% Preámbulo compartido de los manuales K-QGMIP (Fase 8).
\usepackage[utf8]{inputenc}
\usepackage[T1]{fontenc}
\usepackage[spanish,es-tabla,es-noshorthands]{babel} % español, "Tabla", sin shorthands
\usepackage{mathptmx}            % Times New Roman (texto y matemáticas)
\usepackage{microtype}           % mejor justificación (kerning/protrusión)
\usepackage[letterpaper,margin=2.5cm]{geometry}
\usepackage{amsmath,amssymb}
\usepackage{graphicx}
\usepackage{booktabs}
\usepackage{xcolor}
\usepackage{caption}
\usepackage{float}
\usepackage{enumitem}
\usepackage{fancyhdr}
\usepackage{listings}
\usepackage{algorithm}
\usepackage{algpseudocode}
\usepackage[hidelinks,colorlinks=true,linkcolor=black,urlcolor=blue,citecolor=black]{hyperref}

% --- Ajustes de localización (babel ya pone Figura/Tabla/Contenidos) ---
\floatname{algorithm}{Algoritmo}
\captionsetup{labelfont=bf,font=small}
\frenchspacing

% --- Cabeceras/pies ---
\pagestyle{fancy}
\fancyhf{}
\fancyhead[L]{\small K-QGMIP}
\fancyhead[R]{\small Universidad de Caldas — ADA 2026-1}
\fancyfoot[C]{\thepage}
\renewcommand{\headrulewidth}{0.4pt}

% --- Colores y estilo de código ---
\definecolor{codebg}{rgb}{0.97,0.97,0.97}
\definecolor{codekw}{rgb}{0.10,0.10,0.55}
\definecolor{codecm}{rgb}{0.25,0.50,0.30}
\definecolor{codestr}{rgb}{0.60,0.20,0.20}
\lstdefinestyle{py}{
	language=Python,
	backgroundcolor=\color{codebg},
	basicstyle=\ttfamily\footnotesize,
	keywordstyle=\color{codekw}\bfseries,
	commentstyle=\color{codecm}\itshape,
	stringstyle=\color{codestr},
	showstringspaces=false,
	breaklines=true,
	frame=single,
	rulecolor=\color{gray!40},
	numbers=left,
	numberstyle=\tiny\color{gray},
	captionpos=b,
}
\lstdefinestyle{sh}{
	backgroundcolor=\color{codebg},
	basicstyle=\ttfamily\footnotesize,
	breaklines=true,
	frame=single,
	rulecolor=\color{gray!40},
	captionpos=b,
}
\lstset{style=py}
% Permite acentos/UTF-8 dentro de los entornos lstlisting.
\lstset{literate=%
	{á}{{\'a}}1 {é}{{\'e}}1 {í}{{\'i}}1 {ó}{{\'o}}1 {ú}{{\'u}}1
{Á}{{\'A}}1 {É}{{\'E}}1 {Í}{{\'I}}1 {Ó}{{\'O}}1 {Ú}{{\'U}}1
{ñ}{{\~n}}1 {Ñ}{{\~N}}1 {ü}{{\"u}}1 {¿}{{?`}}1 {¡}{{!`}}1}

% Atajos matemáticos del dominio
\newcommand{\dk}{\delta_k}
\newcommand{\EMD}{\mathrm{EMD}}
\newcommand{\bigO}{\mathcal{O}}
\newcommand{\R}{\mathbb{R}}

% Formato de fecha personalizado (mes y año), sin dependencias externas.
\newcommand{\mesano}{%
	\ifcase\month\or enero\or febrero\or marzo\or abril\or mayo\or junio\or
	julio\or agosto\or septiembre\or octubre\or noviembre\or diciembre\fi
	\space de \the\year%
}



\begin{document}

% ====================== PORTADA ======================
\begin{titlepage}
	\centering
	\vspace*{1.5cm}
	{\large\textbf{Universidad de Caldas}}\\[0.2cm]
	{\normalsize Facultad de Inteligencia Artificial e Ingenierías}\\[0.2cm]
	{\normalsize Análisis y Diseño de Algoritmos — 2026-1}\\[2.5cm]
	{\Huge\textbf{Proyecto K-QGMIP}}\\[0.6cm]
	{\LARGE\textbf{Manual de Usuario}}\\[0.4cm]
	{\large Análisis de $k$-particiones de mínima información}\\[0.3cm]
	{\large Estrategias \textbf{KGeoMIP} y \textbf{KQNodes}}\\[3cm]
	{\small\itshape Guía de instalación y uso para usuarios finales}\\[1.5cm]
	{\large\textbf{Autores}}\\[0.2cm]
	{\normalsize Milton Pablo Arias Rincon - 38790}\\[0.1cm]
	{\normalsize Nelson Daniel Estrada Leiton - 9665}\\[1.5cm]
	{\small \mesano}
	\vfill
\end{titlepage}

\tableofcontents
\newpage

% ====================== 1. INTRODUCCIÓN ======================
\section{Introducción y visión general}

\subsection{¿Qué hace el software?}
K-QGMIP analiza un \textbf{sistema} de variables binarias (descrito por una matriz
de transición) y encuentra la mejor manera de \textbf{dividirlo en $k$ partes}
perdiendo la menor cantidad de información posible. En pocas palabras: responde a
la pregunta \emph{``si tuviera que romper este sistema en $k$ pedazos, ¿por dónde
	lo cortaría para que el corte ``duela'' lo menos posible?''}.

\subsection{¿Para qué sirve?}
Es una herramienta de \textbf{Teoría de la Información Integrada} (IIT), usada para
medir cuán ``integrado'' está un sistema (cuánto depende el todo de sus partes).
Aplicaciones típicas: análisis de redes neuronales o genéticas, estudio de
modularidad de sistemas dinámicos, y experimentación académica sobre IIT.

\subsection{Conceptos básicos}
\begin{itemize}
	\item \textbf{Sistema:} conjunto de $n$ nodos que pueden estar en estado $0$ o $1$ (binarios) o probabilidades continuas.
	\item \textbf{$k$-partición:} una división del sistema en exactamente $k$ grupos
	      (bloques). Con $k=2$ es un solo corte; con $k=3,4,5$ son cortes múltiples.
	\item \textbf{Partición de mínima información (MIP):} la división que menos
	      información destruye. El número que mide ese ``daño'' se llama \textbf{pérdida}
	      (símbolo $\varphi$ o $\dk$): \emph{cuanto más pequeño, mejor partición}.
\end{itemize}

\begin{quote}\small
	\textbf{Idea clave:} el programa prueba muchos cortes posibles y se queda con el
	que produce la \emph{pérdida más baja}. Ese es el resultado que usted busca.
\end{quote}

\subsection{Capacidades y limitaciones}
\begin{itemize}
	\item \textbf{Sí puede:} analizar sistemas de hasta $n \approx 25$ nodos con las
	      dos estrategias núcleo KGeoMIP y KQNodes, calculando la pérdida real $\dk$ para
	      cualquier $k\in\{2,3,4,5\}$. La tabla de evaluación oficial está resuelta al
	      100\,\% hasta N25A.
	\item \textbf{Con cuidado:} a $n=25$ el subsistema completo tarda del orden de uno
	      o dos minutos por fila y usa varios GB de RAM (pico $\approx 7{,}9$~GB en KGeoMIP),
	      así que conviene una máquina con 16~GB.
	\item \textbf{No puede:} garantizar la partición óptima global en sistemas
	      grandes (usa heurísticas voraces), ni manejar $n \gg 25$ en un equipo común.
\end{itemize}

% ====================== 2. REQUISITOS ======================
\section{Requisitos del sistema}

\begin{table}[H]
	\centering\small
	\caption{Requisitos mínimos y recomendados.}
	\begin{tabular}{@{}lll@{}}
		\toprule
		                  & Mínimo                      & Recomendado                            \\
		\midrule
		Sistema operativo & Linux / macOS / Windows 10+ & Linux (probado en Arch)                \\
		Python            & 3.14.5                      & 3.14.5                                 \\
		RAM               & 4 GB ($n\le 12$)            & 8 GB ($n\le 20$); 16 GB para $n=25$    \\
		Disco             & 500 MB                      & 2 GB (incluye muestras grandes)        \\
		CPU               & 2 núcleos                   & 4+ núcleos (paralelismo PCD)           \\
		\bottomrule
	\end{tabular}
\end{table}

\noindent\textbf{Software necesario:} el gestor de paquetes
\href{https://github.com/astral-sh/uv}{\texttt{uv}} (instala Python y las
dependencias automáticamente). Las bibliotecas (numpy, scipy, pandas, matplotlib,
etc.) se declaran en \texttt{pyproject.toml} y se instalan con un solo comando.

% ====================== 3. INSTALACIÓN ======================
\section{Instalación paso a paso}

\paragraph{Paso 1 — Obtener el código.} Clonar el repositorio o descargar el
\texttt{.zip}:
\begin{lstlisting}[style=sh]
git clone <URL-del-repositorio> KQGMIP
cd KQGMIP
\end{lstlisting}

\paragraph{Paso 2 — Instalar uv} (si aún no lo tiene):
\begin{lstlisting}[style=sh]
curl -LsSf https://astral.sh/uv/install.sh | sh   # Linux / macOS
# Windows (PowerShell):
# powershell -c "irm https://astral.sh/uv/install.ps1 | iex"
\end{lstlisting}

\paragraph{Paso 3 — Instalar dependencias.} Desde la carpeta del proyecto:
\begin{lstlisting}[style=sh]
uv sync
\end{lstlisting}
Esto crea el entorno virtual y descarga todas las bibliotecas. No es necesario
activar el entorno: cada comando se ejecuta con \texttt{uv run}.

\paragraph{Paso 4 — Verificar la instalación.}
\begin{lstlisting}[style=sh]
uv run pytest -q        # debe terminar con todos los tests en verde
uv run exec.py          # ejecuta un análisis de ejemplo
\end{lstlisting}

\begin{quote}\small
	\textbf{Nota.} Si \texttt{uv run pytest} pasa, el software quedó correctamente
	instalado. El primer \texttt{uv sync} puede tardar varios minutos según la red.
\end{quote}

% ====================== 4. VIDEO ======================
\section{Video tutorial de instalación y uso}

Se incluye un video demostrativo que cubre la instalación desde cero, la
configuración y un ejemplo completo de ejecución e interpretación.

\begin{center}
	\fbox{\parbox{0.8\textwidth}{\centering\vspace{4mm}
			\textbf{Video tutorial:}
			\href{https://youtu.be/5gfKPRqs9HY}{\texttt{https://youtu.be/5gfKPRqs9HY}}\\[4mm]
		}}
\end{center}

\noindent Contenido del video: (1) instalación de \texttt{uv} y
\texttt{uv sync}; (2) ejecución desde la interfaz web y desde la CLI
; (3) interpretación de la salida; (4) generación de figuras.

% ====================== 5. USO BÁSICO ======================
\section{Guía de uso básico}

\subsection{Preparación de los datos de entrada}
El sistema se describe con una \textbf{TPM} (matriz de transición) en un archivo
CSV: \texttt{data/samples/N\{n\}\{página\}.csv}, con $2^n$ filas y $n$ columnas de
$0$/$1$ o probabilidades continuas, en orden \emph{little-endian} (la fila $0$
es el estado $00\dots0$). El programa la carga automáticamente según el número de nodos y la página.

\noindent Si no tiene una TPM, puede generar una desde la \textbf{línea de comandos}
con el script dedicado (no interactivo con \texttt{--yes}):
\begin{lstlisting}[style=sh]
uv run scripts/generate_tpm.py --n 4              # determinista 0/1 (N4?.csv)
uv run scripts/generate_tpm.py --n 6 --continuous # probabilidades continuas
uv run scripts/generate_tpm.py --n 10 --seed 7 --yes
\end{lstlisting}
El archivo se guarda en \texttt{data/samples/} con el siguiente sufijo libre
(\texttt{A}, \texttt{B}, \dots). También puede generarlo \textbf{desde la interfaz
	web} (§\ref{sec:webui}), sin línea de comandos.

\subsection{Ejecución básica}
El análisis individual se lanza con el subcomando \texttt{run} de \texttt{exec.py},
pasando los parámetros por línea de comandos (no hay que editar ningún archivo):

\begin{lstlisting}[style=sh]
uv run exec.py run --net N10A --k 3 --strategy kgeomip
\end{lstlisting}

\noindent Opciones de \texttt{run}:
\begin{itemize}
	\item \texttt{--net} red a analizar, p.\ ej.\ \texttt{N10A} (obligatorio).
	\item \texttt{--k} número de bloques, $\ge 2$ (por defecto 3).
	\item \texttt{--strategy} \texttt{kgeomip}, \texttt{kqnodes}, \texttt{clustering},
	      \texttt{genetic}, \texttt{annealing}, \texttt{tabu} o \texttt{exhaustivek}
	      (por defecto \texttt{kgeomip}).
	\item \texttt{--page} página de la red (A, B, \dots); \texttt{--state} estado
	      inicial binario; \texttt{--method spectral|kmeans} (solo clustering);
	      \texttt{--condition/--purview/--mechanism} para fijar el subsistema a mano.
\end{itemize}

\noindent Sin argumentos, \texttt{uv run exec.py} imprime la ayuda y corre una
demostración de extremo a extremo.

\subsection{Interpretación de resultados}
La salida muestra la estrategia usada, las distribuciones, la $k$-partición
encontrada (bloques \texttt{B1, B2, \dots}) y la \textbf{pérdida mínima}
($\varphi$). Ejemplo real (KGeoMIP, $k=3$, red N10A):

\begin{lstlisting}[style=sh]
KGeoMIP(k=3) fue la estrategia de solucion.
Distancia metrica: distancia-hamming
Notacion indexado: little-endian

Distribucion marginal del Subsistema:
[ 0. 1.0000 0. 0. 1.0000 1.0000 1.0000 0. 0. 0. ]
Distribucion marginal de la Particion:
[ 0. 1.0000 0. 0.4697 1.0000 1.0000 1.0000 0. 0.4727 0. ]

Mejor 3-Particion:
B1: abcdefghij | ABCEFGHJ
B2: ∅ | D
B3: ∅ | I
Perdida minima ( phi ) = 0.9424

Tiempos de ejecucion:
Horas: 0.00 = Minutos: 0.0 = Segundos: 0.0188
\end{lstlisting}

\noindent\textbf{Cómo leerlo:}
\begin{itemize}
	\item \textbf{Bloques B1–B3:} cada bloque es \texttt{presente | futuro}. Las
	      letras minúsculas (izquierda) son nodos de mecanismo/presente ($t$); las
	      mayúsculas (derecha), de purview/efecto ($t{+}1$). Un lado vacío significa que
	      ese bloque no aporta átomos de ese tiempo.
	\item \textbf{Pérdida mínima ($\varphi$):} $0{,}9424$. Es el indicador clave:
	      \emph{más bajo es mejor}. Sirve para comparar estrategias o valores de $k$.
	\item \textbf{Tiempo:} cuánto tardó el análisis.
\end{itemize}

\subsection{Casos de uso típicos}
\begin{enumerate}
	\item \textbf{Una $3$-partición de un sistema pequeño:}
	      \texttt{uv run exec.py run --net N10A --k 3 --strategy kgeomip}.
	\item \textbf{Comparar $k=2,3,4$:} repita cambiando \texttt{--k}; anote la pérdida
	      de cada uno (crece con $k$).
	\item \textbf{Comparar estrategias:} repita con \texttt{--strategy kqnodes} y con
	      \texttt{--strategy clustering} sobre la misma red.
\end{enumerate}

% ====================== 6. OPCIONES AVANZADAS ======================
\section{Opciones y parámetros avanzados}

\subsection{Parámetros de configuración}
\begin{table}[H]
	\centering\small
	\caption{Parámetros ajustables.}
	\begin{tabular}{@{}llll@{}}
		\toprule
		Parámetro      & Bandera             & Valores                                                             & Por defecto          \\
		\midrule
		$k$            & \texttt{--k}        & entero $\ge 2$ (2–5)                                                & 3                    \\
		estrategia     & \texttt{--strategy} & kgeomip, kqnodes, clustering, exhaustivek, genetic, annealing, tabu & kgeomip              \\
		método         & \texttt{--method}   & spectral, kmeans (solo clustering)                                  & spectral             \\
		página         & \texttt{--page}     & A, B, C, \dots                                                      & la de \texttt{--net} \\
		estado inicial & \texttt{--state}    & cadena de $0$/$1$ de longitud $n$                                   & todo 1               \\
		semilla        & application.py      & entero                                                              & 73                   \\
		\bottomrule
	\end{tabular}
\end{table}

\subsection{Modos de operación}
\begin{itemize}
	\item \textbf{Individual} (\texttt{uv run exec.py run --net \dots}): un sistema,
	      con los parámetros por línea de comandos.
	\item \textbf{Por lotes} (\texttt{uv run exec.py batch [datos.xlsx]}): recorre
	      cada hoja \texttt{*-Elementos} del workbook, ejecuta KQNodes y KGeoMIP para
	      $k \in \{2,3,4,5\}$ y escribe la copia de resultados. Es reanudable y nunca
	      modifica la entrada.
	\item \textbf{Ver resultados} (\texttt{uv run exec.py results [resultados.xlsx]}):
	      imprime la tabla en la terminal (estrategia, $k$, pérdida, tiempo, partición).
	\item \textbf{Exacto} (\texttt{--strategy exhaustivek}): \emph{ground truth} por
	      enumeración; úselo solo en sistemas pequeños ($n \le 6$).
\end{itemize}

\subsection{Opciones de salida y visualización}
\begin{lstlisting}[style=sh]
uv run exec.py benchmark               # estrategias x redes x k -> CSV/Excel en data/results/
uv run exec.py benchmark --quick       # solo N10A (rapido)
uv run scripts/make_figures.py         # graficas estaticas (PNG): escalabilidad/perdida
uv run scripts/make_interactive.py     # graficas interactivas (HTML, Plotly)
uv run scripts/make_viz.py --net N4A --k 3 --strategy kgeomip   # dibuja la particion
\end{lstlisting}
Las gráficas \textbf{estáticas} (PNG) quedan en \texttt{data/results/figures/} y las
\textbf{interactivas} (HTML, navegables con \emph{zoom} y \emph{hover}) en
\texttt{data/results/figures/interactive/}.

\subsection{Optimización de rendimiento}
\begin{itemize}
	\item En sistemas pequeños \texttt{kgeomip} es más rápido; en sistemas grandes
	      ($n\ge 22$) la relación se invierte y \texttt{kqnodes} resuelve más rápido
	      (su preparación es más barata que la tabla de costos geométrica). Si el tiempo
	      importa a $n$ grande, prefiera \texttt{kqnodes}; ejecute ambas y tome el mínimo.
	\item Desactive el \emph{profiling} para medir tiempos reales
	      (\texttt{application.disable\_profiling()}).
	\item El paralelismo por procesos se activa en las rutas que lo soportan; evite
	      sobre-suscribir núcleos (el software ya fija \texttt{OMP\_NUM\_THREADS=1} en los
	      trabajadores).
\end{itemize}

% ====================== 6B. INTERFAZ WEB ======================
\section{Interfaz web (sin línea de comandos)}
\label{sec:webui}
Para quien prefiera no usar la línea de comandos, el proyecto incluye una
\textbf{interfaz web} (Streamlit) que cubre todo el flujo desde el navegador.
Streamlit y Plotly son dependencias de núcleo, de modo que \texttt{uv sync} ya las
instala; basta con lanzar la aplicación:
\begin{lstlisting}[style=sh]
uv run streamlit_app.py
\end{lstlisting}
Se abre en \texttt{http://localhost:8501}. La interfaz tiene tres pasos:
\begin{enumerate}
	\item \textbf{Datos.} Elija una muestra existente o \emph{genere una TPM nueva}
	      (nodos~$n$, determinista o continua, semilla) con un botón; el CSV se crea en
	      \texttt{data/samples/}.
	\item \textbf{Estrategia.} Seleccione la estrategia (KGeoMIP, KQNodes,
	      Clustering, Genética, Recocido, Tabú, ExhaustiveK), el número de bloques~$k$ y,
	      si lo necesita, las máscaras de subsistema (condición, \emph{purview},
	      mecanismo).
	\item \textbf{Resultados.} Al pulsar \emph{Ejecutar} verá la pérdida $\dk$, la
	      $k$-partición de mínima información, su distribución marginal en tabla, y el
	      diagrama \textbf{interactivo} de la partición. Si existe la tabla de
	      \emph{benchmark}, se muestra además la gráfica interactiva $\dk$ vs $k$.
\end{enumerate}
\newcommand{\guishot}[2]{%
	\IfFileExists{../../data/.screenshots/#1}{%
		\begin{figure}[H]\centering
			\fbox{\includegraphics[width=0.92\textwidth]{../../data/.screenshots/#1}}
			\caption{#2}\end{figure}%
	}{}%
}

\guishot{generate_tpm_continuos_n6_example.png}{Paso 1: generar una TPM nueva
	desde el navegador (nodos $n$, probabilidades continuas, semilla). La red queda
	guardada en \texttt{data/samples/} (aquí \texttt{N6B.csv}).}

\guishot{ejecucion_basica_k3_n10_kgeomi.png}{Paso 3: resultado de KGeoMIP con $k=3$
sobre N10A. Se ven la pérdida $\dk=0{,}9424$, la $k$-partición de mínima información,
su distribución marginal reconstruida y el diagrama interactivo por bloques.}

\guishot{tabla_de_evaluacion.png}{Tabla de evaluación (\texttt{.xlsx}) en la web:
	carga del workbook estándar, selección de hojas y vista de resultados tipo consola
	(estrategia, $k$, pérdida, tiempo, partición), equivalente a
	\texttt{uv run exec.py results}.}

\guishot{benchmark_n20.png}{Benchmark $\dk$ vs $k$ (N20A) comparando las siete
	estrategias. Abajo, el selector \emph{Redes a incluir en el benchmark} permite elegir
	qué redes de \texttt{data/samples/} se ejecutan y regenerar el benchmark con ellas.
	Clustering espectral queda arriba (pérdida alta), mientras KGeoMIP, KQNodes y Tabú se
	mantienen en el mínimo.}

% ====================== 7. SOLUCIÓN DE PROBLEMAS ======================
\section{Solución de problemas}

\begin{table}[H]
	\centering\small
	\caption{Errores comunes y solución.}
	\begin{tabular}{@{}p{4.4cm}p{4.0cm}p{4.6cm}@{}}
		\toprule
		Síntoma                                          & Causa                                    & Solución                                                       \\
		\midrule
		\texttt{FileNotFoundError: N\{n\}\{page\}.csv}   & No existe la muestra para ese $n$/página & Generar la red o ajustar \texttt{PAGE}/\texttt{initial\_state} \\
		\texttt{ValueError: k must be >= 2}              & $k<2$                                    & Usar $k\ge 2$                                                  \\
		\texttt{ValueError: strict k-partition \dots}    & La partición dejaría un bloque vacío     & Reducir $k$ (no caben tantos bloques no vacíos)                \\
		\texttt{MemoryError} / proceso ``muerto''        & $n$ demasiado grande (RAM)               & Bajar $n$, o usar \texttt{clustering}                          \\
		\texttt{ModuleNotFoundError}                     & Dependencias no instaladas               & Ejecutar \texttt{uv sync}                                      \\
		\texttt{ModuleNotFoundError: streamlit / plotly} & Dependencias no instaladas               & Ejecutar \texttt{uv sync}                                      \\
		Comando \texttt{uv} no encontrado                & \texttt{uv} no instalado/PATH            & Reinstalar \texttt{uv} y reabrir la terminal                   \\
		\bottomrule
	\end{tabular}
\end{table}

\noindent\textbf{Problemas de instalación.} Si \texttt{uv sync} falla, verifique
la conexión y la versión de Python (\texttt{uv python install 3.14}). En Windows,
ejecute la terminal como usuario normal (no se requieren permisos de
administrador).

\noindent\textbf{Datos de entrada problemáticos.} La TPM debe tener exactamente
$2^n$ filas y $n$ columnas con solo $0$/$1$. Las cadenas
\texttt{initial\_state}/\texttt{conditions}/\texttt{purview}/\texttt{mechanism}
deben tener longitud $n$.

% ====================== 8. EJEMPLOS ======================
\section{Ejemplos y tutoriales}

\subsection{Tutorial básico (4 nodos)}
\begin{enumerate}
	\item Ejecutar \texttt{uv run exec.py run -{}-net N4A -{}-k 3 -{}-strategy kgeomip}.
	\item Leer la pérdida y los bloques. Dibujar la partición:
	      \texttt{uv run scripts/make\_viz.py -{}-net N4A -{}-k 3 -{}-strategy kgeomip}.
\end{enumerate}

\subsection{Caso intermedio (10 nodos, varios $k$)}
Con la red N10A, ejecute KGeoMIP para $k=2,3,4$ y compare la pérdida (debe crecer):
$0{,}4727 \to 0{,}9424 \to 1{,}4229$. Repita con \texttt{kqnodes} para confirmar que
obtiene la misma calidad.

\subsection{Ejemplo avanzado (tabla completa)}
\begin{lstlisting}[style=sh]
uv run exec.py batch                   # llena la tabla de evaluacion (k=2..5)
uv run exec.py results                 # ve la tabla en la terminal
uv run scripts/validate.py correctness # valida contra el exacto
uv run scripts/make_figures.py         # produce las graficas
\end{lstlisting}
Los resultados quedan en \texttt{data/results/} (CSV/Excel) y las figuras en
\texttt{data/results/figures/}.

% ====================== 9. REFERENCIA RÁPIDA ======================
\section{Referencia rápida}

\subsection{Comandos principales}
\begin{table}[H]
	\centering\small
	\begin{tabular}{@{}ll@{}}
		\toprule
		Comando                                          & Descripción                           \\
		\midrule
		\texttt{uv sync}                                 & Instalar dependencias                 \\
		\texttt{uv run exec.py run -{}-net N10A -{}-k 3} & Análisis individual                   \\
		\texttt{uv run exec.py batch}                    & Llenar la tabla de evaluación (Excel) \\
		\texttt{uv run exec.py results}                  & Ver la tabla de resultados            \\
		\texttt{uv run exec.py benchmark}                & Benchmark de estrategias              \\
		\texttt{uv run pytest}                           & Ejecutar los tests                    \\
		\texttt{uv run scripts/make\_viz.py}             & Visualizar una partición              \\
		\bottomrule
	\end{tabular}
\end{table}

\subsection{Glosario}
\begin{description}
	\item[TPM] Matriz de probabilidad de transición que describe el sistema.
	\item[$k$-partición] División del sistema en $k$ bloques no vacíos.
	\item[Pérdida ($\varphi$ / $\dk$)] Información destruida por una partición; menor
	      es mejor.
	\item[KGeoMIP] Estrategia geométrica (rápida, recomendada por defecto).
	\item[KQNodes] Estrategia submodular (más robusta; más lenta en $n$ pequeño, más
	      rápida en $n\ge 22$).
	\item[Clustering] Línea base rápida para sistemas grandes (solo propone).
	\item[Purview / Mechanism] Conjuntos de nodos de efecto ($t{+}1$) y mecanismo
	      ($t$) bajo análisis.
\end{description}

\subsection{Formato de archivos}
\begin{itemize}
	\item \textbf{Entrada:} \texttt{data/samples/N\{n\}\{página\}.csv} — $2^n$ filas,
	      $n$ columnas, $0$/$1$, little-endian.
	\item \textbf{Salida:} resumen por consola; tablas en
	      \texttt{data/results/*.csv}/\texttt{*.xlsx}; figuras PNG en
	      \texttt{data/results/figures/}.
\end{itemize}

\end{document}
